% =====================================================================
%  Kapitel 7 – Einschätzung zum Praktikum (~2–3 Seiten)
%  Ich-/Wir-Form erlaubt.
% =====================================================================
\chapter{Unsere Einschätzung zum Praktikum}
\label{kap:einschaetzung}

\section{Studium und Praktikum}
\todoi{Bezug zu Studieninhalten.}

\section{Praktikumsalltag}
\todoi{Ablauf einer typischen Woche, Laborzeiten; neutral: Verfügbarkeit der Drucker im Labor war eingeschränkt, große Teile (je 15\,h+) mussten eingeplant werden.}

\section{Erwartungen}
\todoi{Erwartungen vorher -- was davon eingetreten ist.}

\section{Erfahrungen}

Fachlich haben wir in diesem Praktikum sehr viel gelernt. Den größten Anteil daran hatte die
Fertigung: Wir haben unsere Kenntnisse in SolidWorks deutlich vertieft, gerade die starken
Rundungen des Gehäuses waren anspruchsvoll zu konstruieren. Beim 3D-Druck großer Teile, beim
Kleben, Schleifen und Lackieren haben wir erfahren, wie viel Zeit und Geduld eine saubere
Oberfläche erfordert, etwa beim Warten zwischen den Lackschichten oder beim Schleifen von Hand.
Dabei haben wir auch gemerkt, wie schnell man an solchen Details perfektionistisch wird.
Überrascht hat uns, wie viel sich durch Kleben lösen lässt; viele der zuvor geplanten
Bohrungen wurden am Ende gar nicht gebraucht.

Hinzu kamen Erfahrungen mit Elektronik: Wir haben viel gelötet, mit dem ESP32 gearbeitet und
gelernt, wie die einzelnen Komponenten im Inneren des Roboters zusammenspielen. Auch Sensoren
wie den Thermalsensor oder Velostat kannten wir vorher nicht. Solche Lösungen können wir uns
gut vorstellen, auch in späteren Projekten wieder einzusetzen.

Innerhalb des Teams Sensorik hat die Zusammenarbeit sehr gut funktioniert. Die Abstimmung mit
den anderen Beteiligten lief ebenfalls gut, wenn sie stattfand, und hat uns jedes Mal
weitergebracht. Viele Entscheidungen mussten wir jedoch selbstständig im Team treffen. Das
zeigte sich am deutlichsten, als wir aus eigener Initiative die Konstruktion des gesamten
Gehäuses übernommen haben. Die meisten Probleme traten während der Entwicklung auf, etwa als
die passive Atemerkennung nicht zuverlässig funktionierte oder die Kupferfolie nicht richtig
leitete. In solchen Situationen mussten wir selbst nach einer Lösung suchen, und das gelang
uns mit der Zeit immer besser.
\todoi{Motivation: Was hat euch angetrieben, was war zäh? (ein bis zwei Sätze)}

\section{Persönliches Fazit}
\todoi{Persönliches Fazit.}
